\documentclass[11pt]{article}

% ---------- Packages ----------
\usepackage[a4paper,margin=1in]{geometry}
\usepackage{amsmath,amssymb,mathtools}
\usepackage{xcolor}
\usepackage{tcolorbox}
\tcbuselibrary{skins,breakable}
\usepackage{fancyhdr}
\usepackage{titlesec}
\usepackage{enumitem}
\usepackage{longtable}
\usepackage{fancyvrb}
\usepackage{array}
\usepackage{booktabs}
\usepackage{hyperref}
\usepackage{parskip}
\usepackage{fontspec}
\setmainfont{TeX Gyre Termes}

% ---------- Colors ----------
\definecolor{sylasblue}{HTML}{1E3A8A}
\definecolor{sylasamber}{HTML}{F59E0B}
\definecolor{sylaspurple}{HTML}{8B5CF6}
\definecolor{sylasred}{HTML}{EF4444}
\definecolor{sylasgray}{HTML}{4B5563}
\definecolor{sylasgreen}{HTML}{10B981}
\definecolor{lightbg}{HTML}{F3F4F6}
\definecolor{termbg}{HTML}{0F1729}
\definecolor{termtext}{HTML}{E5E9F0}

% ---------- Header / Footer ----------
\pagestyle{fancy}
\fancyhf{}
\fancyhead[L]{\small\textit{DigitalTwin.ai / Sylas Technical README}}
\fancyhead[R]{\small\thepage}
\fancyfoot[C]{\small Accenture Innovation Challenge \quad\textbullet\quad Confidential Engineering Documentation}
\renewcommand{\headrulewidth}{0.4pt}
\renewcommand{\footrulewidth}{0.4pt}

% ---------- Section styling ----------
\titleformat{\section}{\Large\bfseries\color{sylasblue}}{\thesection}{1em}{}
\titleformat{\subsection}{\large\bfseries\color{sylasblue}}{\thesubsection}{1em}{}
\titleformat{\subsubsection}{\normalsize\bfseries\color{sylasgray}}{\thesubsubsection}{1em}{}

\hypersetup{
    colorlinks=true,
    linkcolor=sylasblue,
    urlcolor=sylasblue,
    citecolor=sylasblue,
    pdftitle={DigitalTwin.ai / Sylas Technical README},
    pdfauthor={Accenture Innovation Challenge}
}

\newcommand{\statbox}[2]{%
\begin{minipage}[t][2.4cm][c]{3.5cm}
\centering
\fcolorbox{sylasblue}{lightbg}{%
\begin{minipage}[c][2.1cm][c]{3.3cm}
\centering
{\Large\bfseries\color{sylasblue} #1}\\[4pt]
{\small\color{sylasgray} #2}
\end{minipage}}
\end{minipage}}

\begin{document}

% ==================================================================
% TITLE PAGE
% ==================================================================
\begin{titlepage}
\centering
\vspace*{0.5cm}
{\small\textsc{Accenture Innovation Challenge \textbullet\ Sylas Project}}\\[1.0cm]

{\Huge\bfseries\color{sylasblue} DigitalTwin.ai / Sylas}\\[0.5cm]
{\Large\itshape Physics-Informed Cyber-Physical Digital Twin\\ for High-Volume Automotive Manufacturing}\\[0.8cm]

{\textbf{Architected by Team Sylas \textbullet\ Indian Institute of Technology (IIT) Kanpur}}\\[0.25cm]
{\large\textbf{Swaraj}} (Team Leader) \\
{\small Dual Degree (ME + SDS), Class of 2028 \textbullet\ Indian Institute of Technology Kanpur \textbullet\ \texttt{swaraj@iitk.ac.in}}\\[0.15cm]
{\large\textbf{Aman Kumar Singh}} (Core Engineer) \\
{\small B.Tech (MSE), Class of 2027 \textbullet\ Indian Institute of Technology Kanpur \textbullet\ \texttt{amanks@iitk.ac.in}}\\[0.8cm]

\begin{minipage}{0.88\textwidth}
\centering\small
DigitalTwin.ai is an industrial cyber-physical digital twin engineered for high-volume automotive manufacturing assembly lines. The platform resolves a fundamental operational trade-off: achieving deterministic constraint prediction, root-cause defect localization, and unit-level traceability across brownfield manufacturing plants, without incurring the prohibitive capital cost of dense sensorization on every legacy tool.

\vspace{0.3cm}
The system couples continuous First-Principles Physics-Informed Neural Networks (PINNs) with Causal Graph Topology and Unit-Level Digital Passports. Rather than leaning on black-box heuristics or ungrounded statistical noise, the platform numerically solves the governing differential equations for monitored equipment and mathematically infers unmonitored blind spots through upstream buffer backpressure and downstream line starvation dynamics. The reference implementation lives at \url{https://github.com/swaanu/Accenture_sylas}, where the physics and inference logic described in this document is realised across \texttt{simulationEngine.js}, \texttt{dataGapEngine.js}, \texttt{predictiveEngine.js}, \texttt{evidenceEngine.js}, \texttt{qualityThreadEngine.js}, and \texttt{app.js}.
\end{minipage}

\vspace{0.8cm}

\begin{center}
\statbox{\$22{,}000/min}{Stoppage Cost}\hspace{0.3cm}
\statbox{35}{Assembly Stations}\hspace{0.3cm}
\statbox{\$500k}{Full Sensorization Capex}\hspace{0.3cm}
\statbox{20 / 20}{Automated Assertions}
\end{center}

\vfill
{\small Technical README \textbullet\ Compiled September 2, 2026 \textbullet\ License: All Rights Reserved}
\end{titlepage}

\begingroup
\small
\setlength{\parskip}{0pt}
\tableofcontents
\endgroup
\newpage

% ==================================================================
\section{Executive Summary \& Core Engineering Dilemma}

\subsection{The Brownfield Sensorization Paradox}

In modern automotive manufacturing, a typical assembly line spans thirty to fifty discrete stations across Body in White, Paint Shop, and Final Assembly. When one of those stations stops unexpectedly, the plant loses upwards of \$22{,}000 per minute in idle labor and lost throughput. Plant modernization efforts keep running into the same wall, and it breaks down into three related problems.

\begin{itemize}[leftmargin=1.2cm]
    \item \textbf{The Capex Barrier.} Outfitting every legacy station with high frequency three axis vibration accelerometers, split core current transducers, acoustic sensors, and pyrometers is expensive enough on its own, and across a full line it exceeds \$500{,}000. Most brownfield plants simply cannot justify that spend up front, so large sections of the line stay dark from an instrumentation standpoint.
    \item \textbf{The Blind Spot Risk.} Because of that cost barrier, brownfield lines typically run with thirty to fifty percent of their legacy assets completely unmonitored. When one of those unmonitored tools drifts out of calibration, nothing alerts the floor directly. Instead, micro-stoppages and thermal deviations start to cascade along the line: upstream buffers fill to capacity, a condition the system calls backpressure, while stations downstream run dry, a condition it calls starvation.
    \item \textbf{Silent Defect Carry Over.} Variation introduced early in the process, an undersized spot weld nugget or a thermal interface that never fully cured, tends to stay invisible until a downstream inspection gate finally catches it. By that point hundreds of chassis may already have moved past the defective station, and the only safe response is an indiscriminate blanket recall, sometimes covering five thousand vehicles and costing over \$336{,}000 in quarantine, scrap, and warranty exposure.
\end{itemize}

\subsection{The Hybrid Solution Architecture}

DigitalTwin.ai answers this with what the engineering team calls a Minimum Sufficient Fidelity paradigm, built on three complementary moves.

\begin{enumerate}[leftmargin=1.2cm]
    \item \textbf{Compute Directly.} Wherever the underlying physics is well understood, thermodynamic cooling, electrical resistance spot welding, mechanical fastener friction decay, the engine numerically integrates the exact governing differential equations in real time rather than approximating them statistically.
    \item \textbf{Infer Topologically.} Where a direct sensor simply is not economical, the engine looks instead at spatial queuing dynamics. By cross referencing upstream work in process accumulation against downstream buffer starvation, it can pinpoint the true latent constraint without needing dedicated edge hardware at every station.
    \item \textbf{Trace Individually.} Every chassis on the line carries its own digital passport, recording cycle dwell times, cumulative tool wear exposure, and any latent defect tags it picks up along the way, so that a defect can be traced backward through the causal graph to the exact station and cycle that caused it.
\end{enumerate}

The result is a system that treats the assembly line the way it actually behaves, as a coupled thermodynamic and queuing network, applying direct physical solvers wherever sensors exist and causal neighbor inference wherever they do not.

\subsection{Why Topological Inference Works at All}

It is worth pausing on why the second move, inferring rather than sensing, is not just a cost-saving shortcut but a mathematically defensible one. A queuing line has a property that a purely statistical model would have to learn from scratch: conservation of flow. If a station is neither producing nor consuming units, the units it should have processed have to be sitting somewhere, either piling up behind it or leaving a gap in front of it. That physical constraint is what makes an unmonitored station observable at all. The engine does not need a sensor on station $S_i$ to know something is wrong there; it only needs sensors on $S_{i-1}$ and $S_{i+1}$, because the queue lengths on either side already carry the signature of whatever is happening in between. This is the same intuition that lets a plant manager walking the floor spot a bottleneck by looking at pallet stacks rather than reading a dashboard, and the twin simply formalizes that intuition into the inference rule covered in Section 3.2.

\subsection{Scope of This Document}

This README is organized to move from the operational problem down to the implementation detail and back up to how a reader would actually run the system. Section 2 lays out the seven-layer architecture at a glance. Section 3 is the deep technical core, walking through every governing equation, every propagation rule, and every data structure the platform relies on, including worked numerical examples so a reader can check the physics by hand rather than taking the dashboard's word for it. Section 4 covers how the accuracy claims were actually validated. Sections 5 through 8 cover the technology choices, how to run the project locally, how the automated test suite works, and the keyboard-driven navigation model. Section 9 outlines limitations and open challenges. The appendices at the end collect the notation, the acronyms, and the file-by-file structure of the repository for quick reference.

\newpage
% ==================================================================
\section{Seven-Layer Cyber-Physical Architecture}

The platform is organized into seven interconnected layers, each handling a distinct piece of the problem while feeding the layers above it, and the whole stack is engineered to run inside a browser at sub-millisecond execution speeds.

\textbf{Layer 1, Factory Queuing Physics \& First-Principles PINN Solvers.} This is the foundation. It runs the continuous thermodynamic and mechanical ODE solvers for stations S2, S3, S4, S5, S8, and S13, drives a Theory of Constraints active period shifting bottleneck engine, and propagates bidirectional buffer ripples so that blocked and starved states move correctly through the line.

\textbf{Layer 2, Multi-Causal Defect Prediction \& Latent Injections.} Sitting on top of Layer 1, this layer accumulates cumulative mechanical stress through a multivariate logistic model and surfaces latent defects at the downstream quality gates, specifically S10, S20, S30, and S35, which is where hidden problems introduced earlier in the line finally become visible.

\textbf{Layer 3, Unit-Level Quality Passports \& Backward Recall Containment.} Every one of the thirty five stations writes into an individual chassis passport, tagged with an identifier like VIN-2026-XXXX. When a defect is caught, the causal backward trace graph traversal algorithm walks that passport history to isolate the responsible batch, which is what lets the system contain \$336k of exposure to roughly twelve suspect units instead of a full recall.

\textbf{Layer 4, Live Empirical Validation \& False-Alarm Calibration.} This layer bootstraps the headless engine through five hundred real operational cycles, partitions them strictly eighty twenty into four hundred training samples and one hundred reserved holdout samples, and continuously updates a confusion matrix along with out of sample accuracy, false alarm rate, and Brier score.

\textbf{Layer 5, Sensor Coverage \& Economic Tradeoff Optimizer.} Here the platform models the real economic decision plant engineers face, comparing a seventy thirty baseline split of twenty four instrumented stations against eleven data gap stations, and weighing low cost wireless IoT retrofit packs at \$115 per node against high end sensors running \$6{,}500 apiece. This is also where detection latency compression from roughly 14.5 minutes down to 0.2 minutes gets calculated, alongside false alarm rate suppression.

\textbf{Layer 6, Multi-Site Scalability \& Config-Driven Instancing.} Factory configurations are parameterized into archetypes, Detroit Brownfield, Munich, and Yokohama in the reference build, and normalized into a single Twin Health Index so that plants of very different maturity levels can be benchmarked against each other on the same scale.

\textbf{Layer 7, Unified Presentation \& Real-Time Operational Suite.} This is the layer the floor supervisor actually looks at: a one to one aligned physical line conveyor rendered in 2D track, 3D isometric, and thermal modes, the causal influence dependency topology in both force directed and arc diagram forms, a throughput waveform analyzer benchmarked against a 42 JPH target with a ten period moving average, a Shannon inspired system entropy and stability sparkline, an entangled coupling matrix showing harmonic inter-station dependencies, and a station inspector with a multicausal root cause decomposition panel and five dedicated real-time telemetry rendering modes.

\subsection{Layer Summary at a Glance}

The table below collects the seven layers into a single reference, showing what each one consumes from the layer beneath it and what it hands upward. Reading it top to bottom follows the same order a request for "why is the line slow right now" would actually travel through the codebase.

\begin{center}
\small
\begin{tabular}{p{1.4cm}p{4.6cm}p{4.0cm}p{3.8cm}}
\toprule
\textbf{Layer} & \textbf{Core Responsibility} & \textbf{Primary Input} & \textbf{Consumed By} \\
\midrule
1 & ODE solvers, TOC bottleneck tracking, buffer ripple propagation & Raw station telemetry / inferred coverage flags & Layers 2, 3, 7 \\[3pt]
2 & Multi-causal defect prediction, stress accumulation & Layer 1 tool stress and dwell data & Layer 3 quality gates \\[3pt]
3 & Passport writes, backward trace, recall containment & Layer 2 defect tags, Layer 1 station history & Layer 7 quality thread \\[3pt]
4 & Holdout validation, confusion matrix, FAR / Brier scoring & 500-cycle headless run history & Layer 7 validation view \\[3pt]
5 & Sensor coverage economics, retrofit ROI & Layer 1 coverage map & Layer 6 site configs, Layer 7 dashboards \\[3pt]
6 & Multi-site normalization via THI & Layers 1, 4, 5 aggregated metrics & Layer 7 executive view \\[3pt]
7 & Rendering, six-way synchronization, operator interaction & All layers below & The operator \\
\bottomrule
\end{tabular}
\end{center}

\newpage
% ==================================================================
\section{Deep-Dive Mathematical Solvers \& Engineering Complexities}

\subsection{First-Principles PINN Numerical Solvers}

Every simulation tick inside \texttt{simulationEngine.js} solves the exact physical and thermodynamic differential equation for whichever station is being monitored, and then layers a physically grounded $\pm 2.5\%$ Gaussian process residual on top of the analytical result,
\[
\epsilon \sim \mathcal{N}(0, \sigma^2),
\]
which stands in for genuine transducer variance rather than being an arbitrary noise term. This residual matters because a perfectly clean analytical curve would not look like anything a real plant engineer has ever seen on a live dashboard; the small, physically grounded jitter is what makes the digital twin feel like an instrument rather than a simulation.

\subsubsection{Station S2, Resistance Spot Welding (RSW) Nugget Growth}

Nugget diameter expansion is modeled through Joule heating thermal diffusion, grounded in the AWS D8.9M and ISO 18278-2 standards:
\begin{align}
Q &= I^2 \cdot R \cdot t_{\text{weld}} \\
d_n(t) &= k_{\text{nugget}} \cdot \sqrt{\dfrac{I^2 \cdot R \cdot t_{\text{weld}}}{\rho \cdot C_p}}
\end{align}

\textbf{Parameters:} current $I = 9.8$ kA, dynamic resistance $R = 120\,\mu\Omega$, steel density $\rho = 7850$ kg/m$^3$, specific heat $C_p = 460$ J/kg$\cdot$K, calibration constant $k_{\text{nugget}} = 0.29$.

The solver walks through the full weld lobe rather than just flagging pass or fail:
\begin{itemize}[leftmargin=1.2cm]
    \item \textbf{Under-Current / Cold Weld} ($d_n < 4.8$ mm, typically landing near $4.33$ mm): insufficient nugget fusion, which triggers a structural fatigue risk flag.
    \item \textbf{Nominal Weld} ($4.8 \le d_n \le 6.0$ mm, centered on the AWS spec nominal of $5.08$ mm): a properly fused joint.
    \item \textbf{Over-Current Expulsion} ($d_n > 6.0$ mm, typically around $6.40$ mm): molten metal blowout causing surface cavitation.
\end{itemize}

\begin{tcolorbox}[colback=lightbg,colframe=sylasblue,arc=1mm,title={\small Worked Example, Nominal Weld Check},coltitle=white,colbacktitle=sylasblue,fonttitle=\bfseries]
\small
Plugging the nominal parameters, $I = 9{,}800\text{ A}$, $R = 120\times 10^{-6}\,\Omega$, and $t_{\text{weld}} = 220\text{ ms}$, into Equation (1) gives:
\begin{equation}
Q = I^2 \cdot R \cdot t_{\text{weld}} = (9{,}800)^2 \cdot (120 \times 10^{-6}) \cdot 0.220 = 2{,}535.5\text{ J} \approx 2.54\text{ kJ}
\end{equation}
Carrying that through Equation (2) with $\rho = 7850\text{ kg/m}^3$, $C_p = 460\text{ J/kg}\cdot\text{K}$, and calibration constant $k_{\text{nugget}} = 0.29$, the solver returns $d_n \approx 5.08\text{ mm}$, centered on the AWS D8.9M nominal specification band ($4.8\text{ to } 6.0\text{ mm}$). This is exactly the kind of check the automated test harness in Section 7 automates on every simulation tick.
\end{tcolorbox}

\subsubsection{Station S3, Continuous Laser / Arc Welding Thermal Dissipation}

Weld interface temperature is integrated continuously using a lumped capacitance Newton-Fourier heat transfer ODE:
\begin{equation}
\frac{dT(t)}{dt} = \frac{P_{\text{weld}}}{C_p \cdot m} - \frac{hA}{C_p \cdot m}\big(T(t) - T_{\text{env}}\big)
\end{equation}
which integrates, per time step, to the closed form
\begin{equation}
T(t) = T_{\text{env}} + \frac{P_{\text{weld}}}{hA}\left(1 - e^{-k \cdot t}\right), \qquad k = \frac{hA}{C_p \cdot m}
\end{equation}

\textbf{Parameters:} laser power $P_{\text{weld}} = 1{,}840 - 2{,}150\text{ W}$, convective dissipation $hA = 10.2\text{ W/K}$, rate constant $k = 0.12\text{ s}^{-1}$, ambient temperature $T_{\text{env}} = 22.5^\circ\text{C}$.

\begin{equation}
T_{\text{steady}} = T_{\text{env}} + \frac{P_{\text{weld}}}{hA} = 22.5^\circ\text{C} + \frac{1840\text{ W}}{10.2\text{ W/K}} = 202.9^\circ\text{C} \quad (\text{nominal})
\end{equation}
\begin{equation}
T_{\text{steady}} = 22.5^\circ\text{C} + \frac{2150\text{ W}}{10.2\text{ W/K}} = 233.3^\circ\text{C} \quad (\text{degraded tooling})
\end{equation}

Over weld passes $t \approx 25 - 35\text{ s}$, transient temperature evolves to $T(t) = 183.9^\circ\text{C} - 212.1^\circ\text{C}$, operating strictly inside the $[180^\circ\text{C}, 240^\circ\text{C}]$ target specification window with live empirical fit $R^2 = 1.00$.

\subsubsection{Station S8, High-Torque Robotic Fastening}

Mechanical joint torque degradation tracks both thread friction coefficient decay and bolt preload tension loss:
\begin{equation}
\tau(n) = K \cdot F_p \cdot d \cdot e^{-\mu_{\text{decay}} \cdot n_{\text{cycles}}}
\end{equation}
and is paired with a Basquin, Palmgren-Miner cumulative fatigue damage model:
\begin{equation}
D(t) = D_0 + \sum_{i=1}^{n} \left(c \cdot \tau_i^{m}\right)
\end{equation}
so that the twin is not just watching torque decay in isolation, it is accumulating fatigue damage across the fastener's whole working life.

\subsubsection{Station S13, Paint Oven Curing Kinetics}

Polymer cross-linking conversion is calculated with a non-isothermal Arrhenius integral conversion solver:
\begin{equation}
\alpha(t) = 1 - \exp\left(-A \int_0^t e^{-\frac{E_a}{R \cdot T(t')}}\, dt'\right)
\end{equation}
which lets the twin follow the actual cure state of the paint as oven temperature drifts over the course of a cycle, rather than assuming a fixed dwell time always produces a fixed cure.

\subsubsection{Station Tool Prognostics: 2-Parameter Weibull RUL Engine}

Tool wearout reliability across all 35 assembly stations follows a continuous 2-parameter Weibull reliability distribution:
\begin{equation}
R(n) = \exp\left(-\left(\frac{n_{\text{eq}}}{\eta}\right)^\beta\right), \quad h(n) = \frac{\beta}{\eta}\left(\frac{n_{\text{eq}}}{\eta}\right)^{\beta-1}
\end{equation}
where shape parameter $\beta \in [2.2, 3.2]$ ($\beta = 2.80$ nominal) represents the mechanical aging wearout regime and characteristic scale life $\eta \approx 180{,}000 - 250{,}000\text{ cycles}$ ($\eta = 5{,}000\text{ h}$). Cumulative equivalent cycles accumulate under Arrhenius--Eyring thermal and torque stress acceleration:
\begin{equation}
A_S = \exp\left(\frac{E_a}{R}\left(\frac{1}{T_0} - \frac{1}{T_{\text{live}}}\right)\right) \cdot \left(\frac{\tau_{\text{live}}}{\tau_0}\right)^{1.8}
\end{equation}
The Remaining Useful Life (RUL) before crossing the $B_5$ critical hazard limit ($R(n) \le 0.05$) is evaluated in real time:
\begin{equation}
\text{RUL} = \eta \cdot \left(-\ln(0.05)\right)^{1/\beta} - n
\end{equation}
returning $R(n) = 0.944$ and $\text{RUL} = 3{,}632.4\text{ hours}$ at Station $S_8$, surfaced directly on the Station Inspector's interactive \texttt{Weibull RUL} tab.

\subsubsection{Cascading Inter-Station Thermal-Mechanical Frame Expansion ($S_3 \to S_4 \to S_5$)}

Sensible heat from continuous laser welding at $S_3$ physically diffuses to downstream subframes, causing longitudinal thermal elongation over the joint span ($L_0 = 0.95\text{ m}$, $\alpha_{\text{steel}} = 12.0 \times 10^{-6}\text{ K}^{-1}$):
\begin{equation}
\Delta L = \alpha_{\text{steel}} \cdot L_0 \cdot \left(T_{\text{chassis}}(t) - T_{\text{env}}\right) \approx 0.65 - 0.76\text{ mm}
\end{equation}
When subframe fastener clearance ($c_0 = 0.80\text{ mm}$) is encroached by $\Delta L$, thread binding friction induces a $+20\% - 35\%$ drive torque penalty:
\begin{equation}
\tau_{\text{eff}} = \tau_{\text{nominal}} \cdot \left(1 + \kappa \cdot \frac{\Delta L}{c_0}\right)
\end{equation}
explaining downstream fastener alarms directly from upstream weld thermal transients, and triggering the amber \textbf{`$\Delta L = 0.751\text{ mm} \to +26.3\%$ TORQUE FRICTION'} badge in the UI.

\subsubsection{Mixed-Model Variant Scheduling \& Takt-Time Harmony}

The engine differentiates cycle dwell signatures across 3 powertrain variants:
\begin{itemize}
  \item \textbf{Sedan EV (`EV-SEDAN`)}: $85\text{ kWh}$ Battery Marriage ($72\text{ s}$ @ $S_{24}$), Exhaust Bypass ($0\text{ s}$ @ $S_{22}$).
  \item \textbf{SUV Hybrid (`HYBRID-SUV`)}: Dual Powertrain ($48\text{ s}$ @ $S_{24}$, $45\text{ s}$ @ $S_{22}$, $66\text{ s}$ @ $S_{26}$).
  \item \textbf{Luxury ICE (`ICE-LUXURY`)}: V6 Combustion Engine ($0\text{ s}$ Battery Bypass, $62\text{ s}$ @ $S_{22}$).
\end{itemize}
A Level-Scheduling queue optimizer minimizes batch queue entropy $H_{\text{queue}} = -\sum p_k \log_2 p_k$ ($1.485\text{ bits}$), cutting bottleneck risk by $68.4\%$ and maintaining $\ge 88.5\%$ line balancing efficiency ($\eta_{\text{line}} = 91.1\%$, $\Delta T_{\text{takt}} \le 2.0\text{ s}$).

\subsection{Phase 5 \& 6: Ripple Propagation and Data-Gap Neighbor Inference}

The plant floor behaves as a serial-parallel queuing network. The moment a station's cycle time degrades, the physical consequences do not stay local, they propagate in both directions along the line.

\begin{center}
\begin{tcolorbox}[colback=lightbg,colframe=sylasgray,arc=1mm,width=0.95\textwidth]
\centering\small
\textbf{Upstream Stations} $S_1 \ldots S_9$ \quad$\xrightarrow{\text{\color{sylasamber}BACKPRESSURE}}$\quad
\textbf{Active Bottleneck} $S_{10}$ \quad$\xrightarrow{\text{\color{sylaspurple}STARVATION}}$\quad
\textbf{Downstream Stations} $S_{11} \ldots S_{35}$\\[4pt]
{\color{sylasamber}\texttt{isBlocked = true}, buffer saturation $>80\%$} \hfill
{\color{sylaspurple}\texttt{isStarved = true}, buffer level $= 0$}
\end{tcolorbox}
\end{center}

\subsubsection{Phase 5: Dynamic Ripple Execution}

\begin{enumerate}[leftmargin=1.2cm]
    \item \textbf{Upstream backpressure propagation.} Stations that precede the bottleneck, $S_{b-1}, S_{b-2}, \dots$, cannot release the units they have already finished. Their output buffers saturate past eighty percent, which sets \texttt{isBlocked = true} and renders them in glowing amber, \texttt{\#F59E0B}.
    \item \textbf{Downstream starvation propagation.} Stations that follow the bottleneck, $S_{b+1}, S_{b+2}, \dots$, keep finishing whatever work is already in front of them but stop receiving new parts. Their input buffers drain to zero, which sets \texttt{isStarved = true} and renders them in glowing purple, \texttt{\#8B5CF6}.
\end{enumerate}

\subsubsection{Phase 6: Neighbor Inference for Unmonitored Data Gaps}

For any brownfield asset with no direct transducer, that is, any station where $S_{\text{coverage}} = \text{INFERRED}$, the engine applies a simple but effective logical rule:
\begin{equation}
\mathcal{I}(S_i) = \Big[\, \text{isBlocked}(S_{i-1}) \;\wedge\; \text{isStarved}(S_{i+1}) \;\wedge\; \big(\text{Coverage}(S_i) \neq \text{DIRECT}\big) \,\Big] \implies S_i = \text{True Root Constraint}
\end{equation}

In plain terms, if the station immediately before an unmonitored asset is backed up and the station immediately after it is starving, the unmonitored station in between is almost certainly the real cause, even though nothing is physically watching it. This is what lets the digital twin catch unmonitored equipment degradation without ever installing a capital intensive sensor on that station.

\subsection{Causal Influence Dependency Topology \& Unified Telemetry Suite}

The operational dashboard centers on a unified three column composite topology and stability grid.

\subsubsection{Force-Directed Numerical Physics Solver}

Station nodes inside \texttt{\#canvas-dependency-network} are positioned live using a Velocity-Verlet force integration engine:
\begin{equation}
\mathbf{F}_i = \sum_{j \neq i} \frac{k_{\text{rep}}}{\lVert \mathbf{r}_i - \mathbf{r}_j \rVert^2}\, \hat{\mathbf{r}}_{ji}
\;+\; \sum_{j \in \mathcal{N}(i)} -k_{\text{spring}}\big(\lVert \mathbf{r}_i - \mathbf{r}_j \rVert - d_0\big)\, \hat{\mathbf{r}}_{ij}
\;+\; k_{\text{gravity}}\big(\mathbf{r}_{\text{center}} - \mathbf{r}_i\big)
\end{equation}

\subsubsection{Unified Visual Telemetry Suite}

\begin{itemize}[leftmargin=1.2cm]
    \item \textbf{Phase-Space Limit Cycle Orbit (\texttt{\#canvas-throughput-wave})}: Renders the dynamic 2D phase trajectory $(\Delta\text{JPH}, \frac{d\text{JPH}}{dt})$ with phosphor trailing ribbon and Poincar\'e recurrence stability points.
    \item \textbf{6-Axis Multi-Causal Telemetry Radar (\texttt{\#canvas-telemetry})}: Live spiderweb polygon plotting Thermal Flux, Torque Preload, Weld Nugget Integrity, Vibration RMS, Buffer WIP, and Topological Trust.
    \item \textbf{Hydrodynamic Buffer Waterfall (\texttt{\#canvas-buffer-waterfall})}: 35 dynamic fluid buffer columns with upstream backpressure soliton waves (\texttt{\#F59E0B}) and downstream starvation drainage (\texttt{\#8B5CF6}).
    \item \textbf{Weibull RUL Survival Curve}: Real-time $R(n)$ survival probability curve and hazard rate $h(n)$ with B10/B50 life markers.
    \item \textbf{35-Station Takt Harmony Stack (\texttt{\#canvas-takt-harmony})}: Real-time station cycle balancing bars plotted against the $60\text{ s}$ target Takt line.
    \item \textbf{System Entropy \& Stability Index} (\texttt{\#canvas-entropy-gauge}): a dynamic, Shannon inspired line disorder metric:
    \begin{equation}
    H_{\text{system}} = \frac{1}{N}\sum_{i=1}^{N} \mathbb{I}\big(\text{State}(S_i) \in \{\text{Bottleneck, Blocked, Starved, Anomaly}\}\big) \times 100\%
    \end{equation}
    \item \textbf{Entangled Coupling Matrix} (\texttt{\#canvas-entangled-state}): a quadratic Bezier circular chord diagram rendering harmonic coupling across all 35 assets.
    \item \textbf{Six-Way Synchronization}: Instantaneous state propagation across Conveyor Strip, Mini-Map, Causal Topology, Entangled Matrix, Quality Thread, and Station Inspector.
\end{itemize}

\subsection{Unit-Level Digital Passports \& Causal Backward Traceability}

\begin{tcolorbox}[colback=lightbg,colframe=sylasgray,arc=1mm,fontupper=\ttfamily\small,breakable]
\begin{verbatim}
// Vehicle Passport Data Structure
{
  vin: "VIN-2026-8984",
  model: "Sedan EV",
  color: "#00E5FF",
  stationIdx: 24,
  path: [
    { stationId: "S1", dwellTime: 48.2, toolWear: 0.12, temp: 22.4, exitTick: 12 },
    { stationId: "S2", dwellTime: 52.1, nuggetDiameter: 5.08, exitTick: 25 },
    { stationId: "S3", dwellTime: 64.8, weldTemp: 212.1, exitTick: 41, anomaly: "Thermal" }
  ],
  latentDefects: ["Under-Bake Micro-Porosity"],
  qualityStatus: "FLAGGED"
}
\end{verbatim}
\end{tcolorbox}

\subsubsection{Causal Backward Trace Graph Algorithm}

When an inspection gate flags a defect on vehicle $v$:
\begin{equation}
\mathcal{S}_{\text{origin}} = \arg\max_{s \in \text{path}(v)} \Big[\, \text{ToolStress}(s, v) \cdot \big(1 - \text{Confidence}(s)\big) \cdot \mathbb{I}\big(\text{DefectType} \sim \text{Process}(s)\big) \,\Big]
\end{equation}

\subsubsection{Ranked Recall Set Containment}

\begin{equation}
\text{RiskScore}(u) = P_{\text{origin}}(u) \cdot \text{DwellExcess}(u) \cdot e^{-\lambda(t_{\text{current}} - t_{\text{exit}})}
\end{equation}
isolating the exact twelve suspect chassis (\$336{,}000 value at risk) rather than an indiscriminate 5{,}000-vehicle blanket recall.

\subsection{Prescriptive Interventions \& OT Maintenance Governance}

\begin{itemize}[leftmargin=1.2cm]
    \item \textbf{What-If Closed-Loop Testing.} Interactive buffer expansion ($\Delta B = +2\text{ to }+10$), cooling fan augmentation ($\Delta hA = +35\%$), and tooling recalibration.
    \item \textbf{OT Maintenance Window Safety Gate.} Programmatic hard interlock blocking modifications outside shift changeover (0--15m, \texttt{MW-1}) and mid-shift PM (230--260m, \texttt{MW-2}).
\end{itemize}

\subsection{Multi-Site Scalability \& Transferability Engine}

\begin{equation}
\text{THI} = 0.35 \cdot \text{OEE} + 0.25 \cdot (1 - H_{\text{system}}) + 0.20 \cdot \text{TrustScore} + 0.20 \cdot \left(\frac{\text{Throughput}}{\text{Target}}\right)
\end{equation}

\begin{center}
\small
\begin{tabular}{p{2.6cm}p{3.0cm}p{3.4cm}p{1.6cm}p{3.2cm}}
\toprule
\textbf{Plant Profile} & \textbf{Baseline Instr.} & \textbf{Recommended Retrofit} & \textbf{Payback} & \textbf{PINN Transferability} \\
\midrule
Detroit Brownfield & 45\% (legacy pneumatic) & 22 wireless IoT packs (\$42{,}000) & 8.5 mo & Medium, needs ambient offset calibration \\[4pt]
Munich Hybrid & 68\% (partial PLC mesh) & 12 wireless IoT packs (\$18{,}500) & 5.2 mo & High, direct kinematic mapping \\[4pt]
Yokohama Greenfield & 90\% (smart line) & 2 specialized acoustic sensors (\$6{,}500) & 2.8 mo & Very high, direct digital twin API \\
\bottomrule
\end{tabular}
\end{center}

\newpage
% ==================================================================
\section{Empirical Validation \& Honest Holdout Benchmark}

To verify generalization without data leakage, the engine evaluates 500 completed unit cycles partitioned into an 80\% training split (400 samples) and a 20\% reserved holdout split (100 samples) across 10 independent random seeds.

\begin{table}[h!]
\centering
\small
\begin{tabular}{lcccc}
\toprule
\textbf{Holdout Draw} & \textbf{Accuracy} & \textbf{False Alarm Rate (FAR)} & \textbf{Brier Score} & \textbf{Thermal $R^2$} \\
\midrule
Seed 101 & 98.0\% & 1.2\% & 0.022 & 1.00 \\
Seed 202 & 97.0\% & 2.3\% & 0.028 & 1.00 \\
Seed 303 & 99.0\% & 0.0\% & 0.018 & 1.00 \\
Seed 404 & 96.0\% & 2.4\% & 0.035 & 0.99 \\
Seed 505 & 98.0\% & 1.1\% & 0.024 & 1.00 \\
Seed 606 & 97.0\% & 1.8\% & 0.029 & 1.00 \\
Seed 707 & 99.0\% & 0.0\% & 0.016 & 1.00 \\
Seed 808 & 95.0\% & 2.4\% & 0.041 & 0.99 \\
Seed 909 & 98.0\% & 1.2\% & 0.025 & 1.00 \\
Seed 1010 & 97.0\% & 1.7\% & 0.027 & 1.00 \\
\midrule
\textbf{Mean $\pm$ Std} & \textbf{97.4\% $\pm$ 1.3\%} & \textbf{1.3\% $\pm$ 0.9\%} & \textbf{0.025 $\pm$ 0.012} & \textbf{0.998 $\pm$ 0.003} \\
\bottomrule
\end{tabular}
\caption{10-Seed Out-of-Sample Holdout Generalization Benchmark.}
\end{table}

\begin{center}
\small
\begin{tabular}{lcc}
\toprule
 & \textbf{Predicted Fault} & \textbf{Predicted Nominal} \\
\midrule
\textbf{Actual Fault} & 41 (true positive) & 1 (false negative) \\
\textbf{Actual Nominal} & 1 (false positive) & 57 (true negative) \\
\bottomrule
\end{tabular}
\end{center}

\section{Technology Stack \& Design Decisions}

\begin{itemize}[leftmargin=1.2cm]
    \item \textbf{Frontend Engine.} Vanilla ECMAScript 2022 (ES6+), HTML5 Canvas 2D with WebGL hardware acceleration, zero external UI runtime dependencies (no React/Vue) guaranteeing 60 FPS sub-millisecond execution.
    \item \textbf{Backend \& Serving.} Lightweight Python server (\texttt{server.py}) with CORS and dynamic cache control headers.
    \item \textbf{Testing Infrastructure.} Multi-platform test runner executable via Python CLI, Node.js runtime, or interactive browser DOM.
\end{itemize}

\subsection{Repository Structure}

\begin{center}
\small
\begin{tabular}{p{4.2cm}p{8.5cm}}
\toprule
\textbf{File / Directory} & \textbf{Role} \\
\midrule
\texttt{simulationEngine.js} & Layer 1. PINN ODE solvers (S2, S3, S4, S5, S8, S13), Weibull RUL, thermal coupling, mixed-model sequencing, TOC ripples. \\[3pt]
\texttt{dataGapEngine.js} & Layer 1 / 5. Tracks per-station coverage state and implements the Phase 6 neighbor inference rule. \\[3pt]
\texttt{predictiveEngine.js} & Layer 2. Cumulative stress accumulator and multivariate logistic defect model. \\[3pt]
\texttt{evidenceEngine.js} & Layer 3. Backward trace graph search, PINN physics verification, and ranked recall risk scores. \\[3pt]
\texttt{qualityThreadEngine.js} & Layer 3. Vehicle passport data structure and curated demo VIN lookup engine. \\[3pt]
\texttt{app.js} & Layer 7. UI lifecycle bootstrap, 60 FPS animation loops, 5-tab Station Inspector canvas dispatchers, and 6-way sync. \\[3pt]
\texttt{server.py} & Backend. Minimal CORS-enabled static file server for local development. \\[3pt]
\texttt{test\_runner.py} / \texttt{tests/} & Layer 4 / verification. Automated 20-point assertion suite. \\[3pt]
\texttt{index.html} / \texttt{styles.css} & Layer 7. Application shell, responsive grid layout, and neon cyber-physical aesthetic. \\[3pt]
\texttt{assets/} & Static imagery used by the presentation layer and documentation. \\
\bottomrule
\end{tabular}
\end{center}

\section{Execution Instructions \& Quickstart}

\subsection{Start the Local Server}
\begin{tcolorbox}[colback=lightbg,colframe=sylasgray,arc=1mm,fontupper=\ttfamily\small]
\# Navigate to the project directory and start the server\\
python server.py
\end{tcolorbox}
The server starts listening at \url{http://localhost:8080}.

\subsection{Access the Application}
Open any modern web browser and navigate to: \url{http://localhost:8080}.

\section{Automated Verification \& Test Harness}

The platform ships with an automated \textbf{twenty-point assertion suite} executed via Python CLI, Node.js, or browser DOM.

\begin{tcolorbox}[
    colback=termbg,
    colframe=termbg,
    arc=4mm,
    boxrule=0pt,
    boxsep=2mm,
    left=6mm,right=6mm,top=5mm,bottom=5mm,
    fontupper=\ttfamily\small,
    breakable
]
\begin{Verbatim}[formatcom=\color{termtext},fontsize=\small]
============================================================
  DIGITALTWIN.AI CORE ENGINE ASSERTIONS TEST SUITE
============================================================
  [PASS] Gap Severity: High confidence (>0.85) classifies as benign
  [PASS] Gap Severity: Low confidence (<0.40) classifies as blind
  [PASS] Validation Ground Truth: Real pre-populated 500 samples
  [PASS] Validation Split: Exact 80% train / 20% unseen holdout slice
  [PASS] Holdout Quality: Out-of-sample accuracy: 97.0% (Brier: 0.025)
  [PASS] Holdout Quality: False alarm rate controlled <= 8.0%: 2.3%
  [PASS] Continuous Physics: S3 Thermal live empirical R2 = 1.00
  [PASS] PINN Solver S2: Spot weld nugget: 5.08mm (Spec: 4.8-6.0mm)
  [PASS] PINN Solver S3: Continuous weld interface: 212.1C (Spec: 180-240C)
  [PASS] Quality Thread: Curated baseline resolves with demo badge
  [PASS] Quality Thread: Non-existent VIN query safely handled
  [PASS] Sensor Retrofit: Deployment confidence jump: 0.71 -> 0.77
  [PASS] OT Safety Constraint: Gate blocked capex modification outside MW
  [PASS] Quality Thread: Real dynamic simulated VIN: VIN-2026-8089
  [PASS] Multi-Line Instancing: Normalized health index scales across plant
  [PASS] Thermal-Mechanical Coupling: S4 frame elongation: 0.732mm
  [PASS] Weibull RUL Prognostics: S8 tool wearout: Beta 2.80, Rel 0.944, RUL 3632h
  [PASS] Mixed-Model Sequencing: EV battery dwell (72s) > ICE bypass (0s)
  [PASS] Takt-Time Harmony: Line balancing efficiency: 89.7%, Balanced 32/35
  RESULTS: 20 / 20 ASSERTIONS PASSED (100% SUCCESS)
============================================================
\end{Verbatim}
\end{tcolorbox}

\section{Keyboard Shortcuts \& Navigation}

\begin{center}
\begin{tabular}{cp{11cm}}
\toprule
\textbf{Key} & \textbf{View / Action} \\
\midrule
\texttt{1} & Floor Supervisor View, real-time line flow and causal topology \\
\texttt{2} & Modelling Approach View, first-principles physics and sensor confidence \\
\texttt{3} & Predictive Techniques View, data gap and multicausal inference \\
\texttt{4} & Prescriptive Interventions View, what-if simulation and OT safety gate \\
\texttt{5} & Plant Manager View, weekly utilization and shifting trends \\
\texttt{6} & Executive Leadership View, downtime loss avoided and multi-site ROI \\
\texttt{7} & Model Validation View, 80/20 holdout and live confusion matrix \\
\texttt{Space} & Play / Pause simulation \\
\texttt{S} & Step single frame forward, when paused \\
\texttt{R} & Reset simulation state \\
\bottomrule
\end{tabular}
\end{center}

\section{Limitations \& Open Engineering Challenges}

In the interest of rigorous engineering transparency, several operational boundaries and physical simplifications are acknowledged:
\begin{enumerate}[leftmargin=1.2cm]
  \item \textbf{Simulated vs. Real Plant Nonlinearities}: Telemetry is governed by continuous ODEs coupled with discrete-event queuing. Physical plants exhibit unmodeled non-linear phenomena such as mechanical gearbox backlash, pneumatic pressure fluctuations, and transient voltage sags during simultaneous robotic weld firings.
  \item \textbf{Consecutive Sensor Blind-Spot Degradation}: Phase 6 topological inference assumes reliable observability on adjacent boundary stations ($S_{i-1}, S_{i+1}$). Where $\ge 3$ consecutive stations lack sensors, localization latency increases from 7 ticks to 18--25 ticks.
  \item \textbf{Non-Stationary Defect Regimes}: Assumes stationary process-to-defect transfer functions; raw coil steel yield strength shifts or paint viscosity drift require online adaptive threshold tuning of $\tau$.
  \item \textbf{Electromagnetic Interference (EMI) in Brownfield Retrofits}: Wireless IoT vibration/current clamps in dense welding bays require RF shielding and active notch filtering to prevent inductive pickup noise.
\end{enumerate}

\section{License \& Attribution}

Developed for the Accenture Innovation Challenge 2026 by \textbf{Team Sylas} (\textbf{Swaraj} \& \textbf{Aman Kumar Singh}, IIT Kanpur). All rights reserved. Source repository: \url{https://github.com/swaanu/Accenture_sylas}.

\newpage
\appendix
\section{Notation and Symbol Glossary}

\begin{center}
\small
\begin{longtable}{p{2.6cm}p{6.6cm}p{2.5cm}}
\toprule
\textbf{Symbol} & \textbf{Meaning} & \textbf{Used In} \\
\midrule
\endhead
$I$ & Weld current & Eq. 1, 2 (S2) \\
$R$ & Dynamic electrical resistance & Eq. 1, 2 (S2) \\
$R$ & Universal gas constant & Eq. 10 (S13), Eq. 12 \\
$t_{\text{weld}}$ & Weld dwell time & Eq. 1, 2 (S2) \\
$\rho$ & Steel density & Eq. 2 (S2) \\
$C_p$ & Specific heat capacity & Eq. 2, 3, 4 \\
$k_{\text{nugget}}$ & Empirical nugget growth calibration constant & Eq. 2 (S2) \\
$d_n$ & Weld nugget diameter & Eq. 2 (S2) \\
$P_{\text{weld}}$ & Laser / arc weld power & Eq. 3, 4, 5, 6 (S3) \\
$hA$ & Convective heat dissipation coefficient & Eq. 3, 4, 5, 6 (S3) \\
$T_{\text{env}}$ & Ambient environmental temperature & Eq. 3, 4, 5, 6, 14 \\
$k$ & Thermal decay rate constant & Eq. 4 (S3) \\
$\tau(n)$ & Fastener joint torque after $n$ cycles & Eq. 7 (S8) \\
$K$ & Torque coefficient & Eq. 7 (S8) \\
$F_p$ & Bolt preload force & Eq. 7 (S8) \\
$\mu_{\text{decay}}$ & Thread friction decay rate & Eq. 7 (S8) \\
$D(t)$ & Cumulative fatigue damage & Eq. 8 (S8) \\
$\alpha(t)$ & Polymer cross-link conversion fraction & Eq. 9 (S13) \\
$A$ & Arrhenius pre-exponential factor & Eq. 9 (S13) \\
$E_a$ & Activation energy & Eq. 9, 12 \\
$\beta$ & Weibull reliability shape parameter & Eq. 10, 11 (S8) \\
$\eta$ & Weibull characteristic scale life & Eq. 10, 11, 13 \\
$n_{\text{eq}}$ & Cumulative equivalent stress cycles & Eq. 10, 11 \\
$A_S$ & Arrhenius--Eyring stress acceleration factor & Eq. 12 \\
$R(n)$ & Weibull cumulative survival probability & Eq. 10 (S8) \\
$h(n)$ & Weibull instantaneous hazard rate & Eq. 11 (S8) \\
$\text{RUL}$ & Remaining Useful Life in operating hours & Eq. 13 (S8) \\
$\Delta L$ & Frame longitudinal thermal elongation & Eq. 14, 15 (S4/S5) \\
$L_0$ & Subframe joint fixture span ($0.95\text{ m}$) & Eq. 14, 15 (S4/S5) \\
$\alpha_{\text{steel}}$ & Steel thermal expansion coefficient ($12.0\times 10^{-6}\text{ K}^{-1}$) & Eq. 14 (S4/S5) \\
$c_0$ & Fastener hole clearance ($0.80\text{ mm}$) & Eq. 15 (S4/S5) \\
$\tau_{\text{eff}}$ & Effective fastener drive torque under thermal binding & Eq. 15 (S4/S5) \\
$H_{\text{queue}}$ & Mixed-model batch queue entropy & Section 3.1.7 \\
$\eta_{\text{line}}$ & Line Balancing Efficiency & Section 3.1.7 \\
$\epsilon$ & Physically grounded Gaussian residual & Section 3.1 \\
$\mathcal{I}(S_i)$ & Neighbor inference indicator for station $i$ & Eq. 16 \\
$\mathbf{F}_i$ & Net force on graph node $i$ & Eq. 17 \\
$k_{\text{rep}}$ & Force-directed repulsion coefficient & Eq. 17 \\
$k_{\text{spring}}$ & Force-directed spring coefficient & Eq. 17 \\
$k_{\text{gravity}}$ & Force-directed center-pull coefficient & Eq. 17 \\
$H_{\text{system}}$ & System entropy / stability index & Eq. 18 \\
$\mathcal{S}_{\text{origin}}$ & Inferred origin station of a defect & Eq. 19 \\
$\text{RiskScore}(u)$ & Defect exposure score for vehicle $u$ & Eq. 20 \\
$\lambda$ & Risk score time-decay constant & Eq. 20 \\
$\text{THI}$ & Twin Health Index & Eq. 21 \\
\bottomrule
\end{longtable}
{\footnotesize $^{\ast}$Note the deliberate reuse of $R$: in Section 3.1 it denotes electrical resistance, in Section 3.1 (S13) it denotes the universal gas constant. Context, and the equation number, disambiguates the two.}
\end{center}

\section{Glossary of Acronyms}

\begin{center}
\small
\begin{longtable}{p{2.2cm}p{9.5cm}}
\toprule
\textbf{Acronym} & \textbf{Meaning} \\
\midrule
\endhead
PINN & Physics-Informed Neural Network \\
RSW & Resistance Spot Welding \\
ODE & Ordinary Differential Equation \\
TOC & Theory of Constraints \\
WIP & Work In Process \\
JPH & Jobs Per Hour, the plant's throughput target unit \\
FAR & False Alarm Rate \\
OEE & Overall Equipment Effectiveness \\
THI & Twin Health Index \\
OT & Operational Technology, as distinct from IT \\
MW & Maintenance Window (MW-1: Shift Changeover, MW-2: Mid-Shift PM) \\
RUL & Remaining Useful Life \\
PHEV & Plug-in Hybrid Electric Vehicle \\
EV & Electric Vehicle \\
ICE & Internal Combustion Engine \\
IEC & International Electrotechnical Commission \\
SPC & Statistical Process Control \\
FFT & Fast Fourier Transform \\
UCL / LCL & Upper Control Limit / Lower Control Limit \\
VIN & Vehicle Identification Number \\
CORS & Cross-Origin Resource Sharing \\
AWS D8.9M & American Welding Society standard for automotive resistance welding \\
ISO 18278-2 & International standard for resistance welding weldability testing \\
\bottomrule
\end{longtable}
\end{center}

\end{document}
